\clearpage
\section*{ТЕРМИНЫ И ОПРЕДЕЛЕНИЯ}
\addtocontents{toc}{\protect\contentsline{section}{ТЕРМИНЫ И ОПРЕДЕЛЕНИЯ}{}{}}

В настоящей выпускной квалификационной работе применяют следующие термины с соответствующими определениями:

\textbf{Биржевая котировка} — опубликованное во времени значение цены финансового инструмента, сформированное в результате биржевых сделок и используемое для анализа рыночной динамики

\textbf{Финансовый временной ряд} — упорядоченная по времени последовательность наблюдений за рыночными переменными, такими как цена открытия, максимум, минимум, цена закрытия и объем торгов

\textbf{OHLCV-данные} — формат представления рыночной информации, включающий цену открытия (\textit{Open}), максимум (\textit{High}), минимум (\textit{Low}), цену закрытия (\textit{Close}) и объем торгов (\textit{Volume})

\textbf{Торговая стратегия} — формализованный набор правил, определяющий условия открытия, удержания и закрытия позиции на финансовом рынке

\textbf{Baseline-стратегия} — простая эталонная стратегия, используемая как нижняя граница качества при сравнении более сложных алгоритмов

\textbf{Признак} — числовая характеристика состояния рынка, используемая моделью для принятия решения или прогнозирования

\textbf{Целевая переменная} — величина, которую модель должна предсказать, например направление изменения цены или доходность следующего периода

\textbf{Обучение с подкреплением} — парадигма машинного обучения, при которой агент последовательно взаимодействует со средой, выбирает действия и получает награду за последствия этих действий

\textbf{Агент} — обучаемая система, принимающая решения в среде на основе наблюдаемого состояния

\textbf{Состояние среды} — совокупность признаков, описывающих доступную агенту информацию о рынке и текущем состоянии портфеля

\textbf{Функция вознаграждения} — правило количественной оценки результата действия агента, используемое для настройки его поведения

\textbf{Equity curve} — кривая стоимости портфеля во времени, отражающая изменение капитала при выполнении стратегии

\textbf{Просадка} — относительное снижение стоимости портфеля от локального максимума до последующего минимума
